\section{The full original Cartier fibre projection}

\begin{theorem}[An actual principal Cartier chart projects its whole original fibre]
\label{thm:intrinsic-numerical-cartier-principal-full-fibre}
\uses{scheme:actual-original-flat-global-section-tensor,def:intrinsic-numerical-canonical-function-field-degree,def:fc-curve-cartier-original-equations,def:fc-native-weighted-regular-section-order,def:fc-finite-fibre-original-coefficient-square,scheme:actual-proper-global-section-valuation-membership,scheme:genuine-ideal-pullback-stalk,lem:cartier-principal-equation-regular,lem:nonflat-dominant-cartier-comap,lem:actual-integral-projection-dimension-bound,thm:fc-curve-cartier-native-local-orders,thm:fc-finite-fibre-all-native-principal-orders,thm:intrinsic-numerical-spec-principal-cartier-weight,thm:intrinsic-numerical-cartier-affine-coefficient-rank}
\lean{EnvelopeProperty.NumericalCartierCurveOrders.finsum_fibre_weightedLocalOrder_principal_spec}
\leanok
Let $A$ be an original Noetherian domain of dimension at most one over
$k$, let $f:Y\to\operatorname{Spec}A$ be an actual finite dominant map
with $Y$ integral, and let $x\ne0$ be an original prime whose actual
localization residue field is finite over $k$. Let $I$ be an actual
Cartier ideal generated globally by the section corresponding to an
original nonzero $a\in A$. Then
\[
 \sum_{y\in f^{-1}(x)}\operatorname{ord}_y(f^{-1}I)
       \operatorname{rk}_k\kappa(y)
 =N(f)\operatorname{ord}_x(I)\operatorname{rk}_k\kappa(x).
\]
All coefficients, ideals, points and the function-field degree are
original. There is no flatness, normal-target or order/residue matching
input. The actual residue-finiteness hypothesis is local geometry;
the later projective-curve consumer must derive it internally from the
original Cartier support.
\end{theorem}
\begin{proof}
The actual finite map makes $Y$ affine and locally Noetherian and,
by the existing integral-morphism dimension bound, gives $\dim Y\leq1$.
Actual dominance from the integral source makes the original comap
Cartier. The native affine ideal-comap formula identifies its whole
global generator with the actual section
$\operatorname{coefficientMap}(f)(a)$. Genuine Cartier regularity
makes this same section nonzero. The original local-equation theorem
therefore identifies every term of its native Cartier order with
the native weighted regular-section order of this same section, for
the same original structural map.

Apply the already proved entire finite-fibre principal-order formula.
Its actual whole coefficient-algebra rank is the same canonical
$N(f)$ by the frozen affine-rank theorem. The frozen native spectrum
principal-Cartier-weight theorem identifies its actual base local order
and residue rank with the original weighted Cartier order. Substitution
proves the whole-fibre identity without a supplied matching equality.
Only the unconditional original coefficient-algebra action from
$\operatorname{sectionsAlgebra}$ is used; no flat base-change theorem
or flatness premise is invoked.
\end{proof}

\begin{theorem}[The whole original fibre projects the actual Cartier order]
\label{thm:intrinsic-numerical-cartier-order-full-fibre}
\uses{def:fc-finite-closed-structure-sections,scheme:actual-original-flat-global-section-tensor,def:intrinsic-numerical-canonical-function-field-degree,def:intrinsic-numerical-spec-native-residue-algebra,thm:fc-scheme-finite-principal-generic-rank,def:fc-curve-cartier-original-equations,def:fc-affine-scheme-local-coefficients,thm:fc-curve-proper-support-points,scheme:genuine-ideal-pullback-stalk,thm:fc-stalk-map-original-order-residue,thm:intrinsic-numerical-cartier-principal-full-fibre,thm:intrinsic-numerical-cartier-order-open-pullback,thm:intrinsic-numerical-cartier-affine-coefficient-rank,thm:intrinsic-numerical-cartier-support-residue-finite,thm:intrinsic-numerical-function-field-degree-affine-rank,thm:fc-curve-cartier-native-local-orders,thm:fc-curve-cartier-ideal-finite,lem:nonflat-dominant-cartier-comap,scheme:actual-proper-global-section-valuation-membership}
\lean{EnvelopeProperty.NumericalCartierCurveOrders.finsum_fibre_weightedLocalOrder}
\leanok
Let $C$ be an original integral projective scheme of dimension at most
one over $k$, and let $g:D\to C$ be an actual finite dominant morphism with
$D$ integral. For every original effective Cartier ideal $I$ and every
original $x\in C$,
\[
 \sum_{y\in g^{-1}(x)}
   \operatorname{ord}_y(g^{-1}I)[\kappa(y):k]
 =N(g)\operatorname{ord}_x(I)[\kappa(x):k].
\]
Here every weight uses the same original structural map, and $N(g)$ is
the actual canonical function-field degree. There is no flatness,
normal-target or matching-order hypothesis.
\end{theorem}
\begin{proof}
Actual dominance from the integral source preserves the Cartier
property of the comap. The original integral dimension bound gives dimension at most one
on each actual full inverse affine chart. Finiteness also makes the
original source locally Noetherian. If $x$ lies outside the actual
support of $I$, its native Cartier order is zero. The actual support
comap formula puts every point above $x$ outside the comap support too.
All terms in the same entire fibre, as well as the original term, vanish.

For $x$ on the support, choose an original Cartier affine chart with
regular equation $c$, and take its actual full inverse affine chart.
Restriction gives a bijection between the same entire original fibre
and the fibre of this actual restricted map: the open projection is
injective, and every original point above $x$ lies in that inverse
chart. The preceding open-order theorem preserves both the original
Cartier orders and their original residue coefficients.

The actual Cartier closed subscheme is finite over $k$. Its actual
global section algebra is finite over $k$, hence Artinian; localization
and passage to its residue quotient remain finite over $k$. The actual
closed-immersion and open-stalk residue algebra isomorphisms therefore
prove finiteness of the original local residue extension needed by the
finite-fibre principal-order theorem. They also identify its local
residue degree with the same $[\kappa(x):k]$.

The original affine principal-order theorem now applies to the actual
restricted map and $c$. Its entire fibre sum uses the actual pulled
section. The original ideal comap and germ identities identify that
section with a generator of the same restricted comap Cartier ideal;
the local-equation theorem identifies its native order with that
Cartier order. Canonical affine global-section isomorphisms preserve
the original coefficient maps and their module ranks. The established
original affine-rank comparison thus changes the multiplier into the
same $N(g)$, while the original stalk-localization isomorphism changes
the local base principal order into $\operatorname{ord}_x(I)$.
Every point, local equation and residue-field comparison in this
calculation is derived from the same original maps. Substituting these
identities in the principal fibre sum proves the formula.
\end{proof}

The global whole-line degree formula and proper-surjective full real nef
reflection and injectivity require subsequent finite-support and
Cartier-ratio assembly. No numerical metric or norm assertion is made
here.
